\subsection{关于两个测度乘积的积分}

在本号中，\(\lambda\) 和 \(\mu\) 分别表示 S 和 T 上的测度，\(\nu\) 表示乘积测度 \(\lambda \otimes \mu\)，定义在 \(S \times T\) 上。此外，若 \(f\) 是 \(S \times T\) 上的正函数，则对于每个 \(s \in S\)，以 \(f_s\) 表示 T 上的函数 \(t \mapsto f(s,t)\)，以 \(I_f\) 表示 S 上的函数 \(s \mapsto \mu^\bullet(f_s)\)。

引理 2. --- 设 \(f\) 是 \(\nu\) -可测的正函数，定义在 \(S \times T\) 上；对于每个紧子集 \(L\)（属于 \(T\)），令 \(I_f^L\) 为函数 \(s \mapsto \mu^\bullet(f_s \varphi_L)\)，定义在 \(S\) 上。则函数 \(I_f^L\) 是 \(\lambda\) -可测的，并且

\[
\mathrm{I} _ {f} = \sup _ {\mathrm{L}} \mathrm{I} _ {f} ^ {\mathrm{L}}\tag{5}
\]

\[
\nu^ {\bullet} (f) = \sup _ {L} \lambda^ {\bullet} (I _ {f} ^ {L}),\tag{6}
\]

其中 L 遍历 T 的紧子集所成的集合。

首先注意到包含关系 \(L \subset L'\) 蕴含 \(I_{f}^{L} \leqslant I_{f}^{L'}\)；另一方面，\(\mathrm{I}_{f}^{\mathrm{L}}(s) = \mu_{\mathrm{L}}^{\bullet}((f_{s})_{\mathrm{L}})\) 对所有 \(s \in S\) 成立。因此，公式 (5) 立即由 §1, No.~2 中给出的负荷 \(\mu^{\bullet}\) 的定义得到。若 K 是 S 的紧子集，L 是 T 的紧子集，则按构造有 \(\nu_{\mathrm{K} \times \mathrm{L}} = \lambda_{\mathrm{K}} \otimes \mu_{\mathrm{L}}\)，而 Ch. V, §8, No.~3 的 Prop. 7 蕴含关系

\[
\nu^ {\bullet} (f \varphi_ {\mathrm{K} \times \mathrm{L}}) = \lambda_ {\mathrm{K}} ^ {\bullet} \big ((\mathrm{I} _ {f} ^ {\mathrm{L}}) _ {\mathrm{K}} \big).\tag{7}
\]

此外，\(S \times T\) 的每个紧子集都包含于某个形如 \(K \times L\) 的紧集；在前一个公式中对所有 K 和 L 取上包络，得到

\[
\nu^ {\bullet} (f) = \sup _ {L} \sup _ {K} \lambda_ {K} ^ {\bullet} \big ((I _ {f} ^ {L}) _ {K} \big) = \sup _ {L} \lambda^ {\bullet} (I _ {f} ^ {L}),\tag{8}
\]

即得 (6)。

最后，Ch. V, §8, No.~3 的 Prop. 7 蕴含 \(I_{f}^{L}\) 在 S 的每个紧子集 K 上的限制是 \(\lambda_{K}\) -可测的；这等价于说 \(I_{f}^{L}\) 是 \(\lambda\) -可测的。

命题 11. --- 设 \(f\) 是 \(\geqslant 0\) 的下半连续函数，定义在 \(X = S \times T\) 上。

\begin{enumerate}
\def\labelenumi{\alph{enumi})}
\item
  函数 \(f_{s}: t \mapsto f(s, t)\) 在 \(\mathbf{T}\) 上下半连续，对每个 \(s \in S\) 均如此。
\item
  函数 \(I_{f}: s \mapsto \int^{\bullet} f(s, t) d\mu(t)\) 在 S 上下半连续，并且
\end{enumerate}

\[
\iint_ {\mathrm{X}} ^ {\bullet} f (s, t) d \nu (s, t) = \int_ {\mathrm{S}} ^ {\bullet} d \lambda (s) \int_ {\mathrm{T}} ^ {\bullet} f (s, t) d \mu (t).\tag{9}
\]

性质 \(a\) 显然，因为映射 \(t \mapsto f(s,t)\) 从 T 到 \(\overline{\mathbf{R}}\) 是 f 与连续映射 \(f\) \(t \mapsto (s,t)\) 从 T 到 X 的复合。为证明 \(b\)，我们将使用如下引理：

引理 3. --- 设 X 是拓扑空间（不论是否 Hausdorff），f 是定义在 X 上的 \(\geqslant 0\) 下半连续函数；则 f 是 X 上下半连续函数的一个递增序列 \((f_{n})_{n\in\mathbf{N}}\) 的极限，并且其中每个函数 \(f_{n}\) 都是开集特征函数的正系数线性组合。

给定两个整数 \(k \geqslant 1\) 和 \(n \geqslant 1\)，令 \(J_{kn}\) 表示区间 \(]k/2^n, +\infty]\)（在 \(\overline{\mathbf{R}}\) 中）的特征函数。对于每个 \(x \in \overline{\mathbf{R}}_+\)，置 \(u_n(x) = 2^{-n} \sum_{k=1}^{n \cdot 2^n} J_{kn}(x)\)；立即可见，序列 \((u_n(x))_{n \geqslant 1}\) 递增并以 \(x\) 为极限。因此函数序列 \(f_n = u_n \circ f\) 递增并收敛于 \(f\)，并且有 \(f_{n} = 2^{-n}\sum_{k = 1}^{n\cdot 2^{n}}\varphi_{\mathrm{U}(k,n)}\)，

其中 \(\mathrm{U}(k,n)\) 是 X 的开集 \(\overline{f} (]k / 2^n, + \infty ])\)。

现在证明 \(b\)。函数 \(I_f\) 是递增有向函数族 \(I_f^L\) 的上包络，其中 \(L\) 遍历 \(T\) 的紧子集所成集合（引理 2），所以只需证明函数 \(I_f^L\) 是下半连续的；然后对 \(L\) 取上包络，即可由 (6) 得到公式 (9)（§1, No.~6, Prop. 5）。

于是令 \(\mathcal{H}\) 为正下半连续函数 \(f\)（定义在 \(\mathrm{S} \times \mathrm{T}\) 上）所成的集合，且满足：\(\mathrm{I}_f^{\mathrm{L}}\) 对于每个紧子集 \(\mathrm{L}\)（属于 \(\mathrm{T}\)）都是下半连续的。根据 §1, No.~6 的 Prop. 5，\(\mathcal{H}\) 中每个递增有向集的上确界属于 \(\mathcal{H}\)。由引理 3，因此只需证明开集 \(\mathrm{W}\)（它是 \(\mathrm{S} \times \mathrm{T}\) 的开集）的特征函数属于 \(\mathcal{H}\)。此外，根据 \(\mathrm{S} \times \mathrm{T}\) 上乘积拓扑的定义，开集 \(\mathrm{W}\) 是一个递增有向族 \((\mathrm{W}_{\alpha})_{\alpha \in \mathrm{A}}\) 的并，这些开集形如

\[
\mathrm{W} = \bigcup_ {1 \leqslant i \leqslant n} \left(\mathrm{U} _ {i} \times \mathrm{V} _ {i}\right),
\]

其中 \(U_{i}\) 在 S 中是开集，\(V_{i}\) 在 T 中是开集；根据上面的注记，只需证明这种开集的特征函数属于 H。令 \(s \in S\)，令 U 为由包含 s 的开集 \(U_{i}\) 所成族的交（该族可能为空）；立即可见，\(\varphi_{\mathrm{W}}(s,t) \leqslant \varphi_{\mathrm{W}}(s',t)\) 对所有 \(s' \in U\) 和 \(t \in T\) 成立，从而通过积分，\(\mathrm{I}_{\varphi_{\mathrm{W}}}^{\mathrm{L}}(s) \leqslant \mathrm{I}_{\varphi_{\mathrm{W}}}^{\mathrm{L}}(s')\) 对所有 \(s' \in U\) 成立。因此 \(I_{\varphi_{w}}^{L}\) 下半连续，命题得证。

推论 1. --- 设 \(f\) 是定义在 \(\mathbf{X} = \mathbf{S} \times \mathbf{T}\) 上的正数值函数；则

\[
\iint_ {\mathrm{X}} ^ {*} f (s, t) d \nu (s, t) \geqslant \int_ {\mathrm{S}} ^ {*} d \lambda (s) \int_ {\mathrm{T}} ^ {*} f (s, t) d \mu (t).\tag{10}
\]

事实上，令 \(g\) 是 \(\mathbf{X}\) 上满足 \(g \geqslant f\) 的下半连续函数；根据 Prop. 11，

\[
\begin{array}{r l} \iint_ {\mathrm{X}} ^ {*} g (s, t) d \nu (s, t) & = \iint^ {\bullet} g (s, t) d \nu (s, t) = \int^ {\bullet} d \lambda (s) \int^ {\bullet} g (s, t) d \mu (t) \\ & = \int^ {*} d \lambda (s) \int^ {*} g (s, t) d \mu (t) \geqslant \int^ {*} d \lambda (s) \int^ {*} f (s, t) d \mu (t). \end{array}
\]

对 \(g\) 取下包络，即得不等式 (10)。

推论 2. --- 设 f 是定义在 S × T 上的数值函数且是 ν-可忽略的。则函数 \(f_{s}: t \mapsto f(s, t)\) 对于 λ-几乎每个 \(s \in S\) 都是 μ-可忽略的。

命题 12. --- 设 f 是定义在 X = S × T 上的 ν-可测正函数。假设 f 是 ν-适度的（相应地，μ 是适度的）。则：

\begin{enumerate}
\def\labelenumi{\alph{enumi})}
\item
  集合 N 由那些 \(s \in S\) 组成，使函数 \(f_{s}: t \mapsto f(s, t)\) 不是 \(\mu\) -可测；对于 \(\lambda\) 而言它是可忽略的（相应地，局部可忽略的）。
\item
  映射 \(s \mapsto \int^{\bullet} f(s, t) \, d\mu(t)\) 是 \(\lambda\) -可测的，并且
\end{enumerate}

\[
\iint_ {\mathrm{X}} ^ {\bullet} f (s, t) d \nu (s, t) = \int_ {\mathrm{S}} ^ {\bullet} d \lambda (s) \int_ {\mathrm{T}} ^ {\bullet} f (s, t) d \mu (t).\tag{11}
\]

我们先证明 \(b)\) 在 \(f\) 是 \(\nu\) -适度的情形。根据引理 2，当存在紧子集 \(L\)（属于 \(T\)）使得 \(f\) 在 \(S \times L\) 外为零时，结论的这一部分成立；因为此时有 \(I_f = I_f^{L'}\)；对于每个紧子集 \(L'\)（属于 \(T\) 且包含 \(L\)），公式 (11) 化为 (6)。特别地，\(b)\) 对于函数 \(f\)（该函数在 \(S \times T\) 的某个紧子集外为零）已得证。另一方面，\(b)\) 在 \(f\) 是 \(\nu\) -可忽略的情形成立（Prop. 11 的 Cor. 1）。由于每个 \(\nu\) -适度函数都是一个 \(\nu\) -可忽略函数与一列紧支集函数之和（§1, No.~9, Cor. 3 of Prop. 14），断言 \(b)\) 在 \(f\) 是 \(\nu\) -适度的情况下成立。

同样，\(b)\) 在 \(\mu\) 集中于紧子集 \(L\)（属于 \(T\)）时显然成立（引理 2）。假设 \(\mu\) 是适度的；则存在具有紧支集的测度序列 \((\mu_n)_{n \in \mathbb{N}}\)，定义在 \(T\) 上，使得有关关系 \(\mu = \sum_{n} \mu_n\)。\(\nu = \sum_{n} \lambda \otimes \mu_n\) 对每个测度 \(b)\) 成立，并且 \(\nu_n = \lambda \otimes \mu_n\)；于是 \(\nu = \sum_{n} \nu_n\)。

我们来证明 \(a\)；记 N 为使 \(s \in S\) 中的 \(f_s\) 不是 \(\mu\) -可测的那些点所成的集合；对于 T 的每个紧子集 L，类似地记 \(\mathbf{N}_{\mathrm{L}}\) 为使 \(s \in S\) 中的 \(f_s \varphi_{\mathrm{L}}\) 不是 \(\mu\) -可测的那些点所成的集合。若 K 和 L 分别是 S 和 T 中的紧集，则 \(f_{\mathrm{K} \times \mathrm{L}}\) 关于测度 \(\nu_{\mathrm{K} \times \mathrm{L}} = \lambda_{\mathrm{K}} \otimes \mu_{\mathrm{L}}\) 是可测的，而 Ch. V, §8, No.~2 的 Prop. 2 表明集合 \(\mathbf{N}_{\mathrm{L}}\) 对 \(\lambda_{\mathrm{K}}\) 局部可忽略；由于 K 任意，\(\mathbf{N}_{\mathrm{L}}\) 对 \(\lambda\) 局部可忽略。

假设 f 在形如 \(K \times L\) 的紧集外为零；则 \(N = N_{L}\)，且 N 包含于 K 中；因此 N 是 \(\lambda\) -可忽略的。同样，如果 f 是 \(\nu\) -可忽略的，Prop. 11 的 Cor. 2 蕴含 N 是 \(\lambda\) -可忽略的。然后结合前面两种情形，即可如上处理 f 是 \(\nu\) -适度的情形。

假设 \(\mu\) 集中于 T 的紧子集 L；则仍有 \(N = N_{L}\)，因而 N 对 \(\lambda\) 局部可忽略。由于每个适度测度都是具有紧支集的测度序列之和（\(\S1\)，No.~9, Cor. 5 of Prop. 14），这个结果立即推广到 \(\mu\) 是适度的情形（利用 \(\S1\)，No.~7 的 Prop. 8）。

注记。--- 令 \((\mathrm{K}_{\alpha})_{\alpha\in\mathrm{A}}\) 是 S 关于 \(\lambda\) 的一个分割，并令 \(M = S - \bigcup_{\alpha\in A} K_{\alpha}\)；类似地，定义 \((\mathrm{L}_{\beta})_{\beta\in\mathrm{B}}\) 和 N（针对 T 上的测度 \(\mu\)）。我们把 \(S'\) 记作 S 的子空间 \(K_{\alpha}\) 与离散空间 M 的和所成的局部紧空间；空间 \(T'\) 类似定义，并置 \(X' = S' \times T'\)。局部紧空间 \(X'\) 是 X 的紧子空间族 \((\mathrm{K}_{\alpha} \times \mathrm{L}_{\beta})_{(\alpha,\beta)\in A\times B}\) 与子空间 \(P = (M \times T) \cup (S \times N)\) 的和；后者是 X 的局部 \(\nu\) -可忽略子集（注意 P 一般不是离散空间）。我们在 §1, No.~8 的 Scholium 中看到，存在测度 \(\lambda'\)，定义在 \(S'\) 上，使得关于 \(\lambda\) 和 \(\lambda'\)，可测函数、正函数的本质上积分、本质可积函数及其积分都相同。按照所引用的 Scholium，将测度 \(\mu'\) 定义在 \(T'\) 上，并与 \(\mu\) 关联；将测度 \(\nu'\) 定义在 \(X'\) 上，并与 \(\nu\) 关联；立即可见，\(\nu'\bullet(f \otimes g) = \lambda'\bullet(f)\mu'\bullet(g)\) 对 \(f \in \mathcal{F}_{+}(S)\) 和 \(g \in \mathcal{F}_{+}(T)\) 成立；因此根据 No.~5 的 Prop. 10，有 \(\nu' = \lambda' \otimes \mu'\)。由于 \(X'\) 的拓扑比 X 的拓扑精细，每个 \(\nu\) -适度函数都是 \(\nu'\) -适度的。这个过程使 Lebesgue--Fubini 定理（Ch. V, §8, No.~4, Th. 1）无需新的证明即可推广到当前情形。

